\section{Discussion}

The audit distinguishes a useful diagnostic question from a deployable
reliability claim. A reference-dependent score asks whether the prediction on a
coarse synthetic view disagrees with a finer image that the benchmark retains.
This can diagnose resolution sensitivity when both views genuinely exist. It
does not establish that a system receiving only the coarse image can reproduce
the same ranking or threshold-transfer behaviour.

Information matching changes the interpretation of baseline comparisons.
Confidence, energy, nearest-neighbour distance, virtual-logit matching, and
received-image consistency can all be computed from the received image plus
source data. Their relative performance therefore concerns the score rather
than hidden input access. A reference-dependent score is retained as an
explicit upper-information diagnostic, not presented as an operational
competitor.

The prospective replication strengthens the scope of this conclusion. Positive
anchor-matched gaps occur for ResNet-18 and MobileNetV3-Small under bicubic,
bilinear, nearest-neighbour and box resampling, with all 36 seed-level
intervals excluding zero. The mechanism is therefore robust to the tested
architecture and operator choices. It remains a directional
fine-versus-coarse comparison, not a causal estimate of information access.

The reveal/mask experiment addresses the narrower identification question.
When scale set, anchor, score and marginal fine predictions are fixed, breaking
only image correspondence reduces error AUROC in all 24 combinations.
Corresponding fine-reference information therefore drives the aligned
diagnostic within this intervention. This does not identify the entire
oracle--received gap, because deployment still substitutes coarser views rather
than permuted finer views.

The data audits have a related purpose. Exact duplicate checks target literal
reuse across splits, whereas guarded spatial blocks target shared pixels and
local spatial dependence between fixed crops. Neither control alone establishes
geographical generalisation. The held-out site and event design therefore
complements, rather than replaces, the lower-level leakage checks.

The multi-event result illustrates why pooling is insufficient. Received-image
consistency had a positive pooled satellite AUROC difference because the two
Hurricane Ian acquisitions favoured it, while both Hurricane Michael
acquisitions favoured confidence. Cluster intervals consequently included zero despite three
positive pooled seed differences. Likewise, mean threshold-transfer coverage
and risk favoured consistency, so preregistered W13 passed, but one seed
reversed the risk ordering and the post-hoc all-seed sensitivity failed.
Consequently, W12 does not support general ranking and W13 does not establish
seed-uniform robustness; the measured-GSD evidence supports event-associated
heterogeneity and audit practice rather than a winning gate.

The held-out Hurricane Idalia UAS/crewed comparison adds a third event, a
different aerial platform and substantially closer physical footprints than the
UAS/satellite pairs. Received consistency ranks crewed-aircraft errors above
confidence for all seeds with positive cluster intervals. This broadens the
real-world result, but the extreme undamaged majority and one Idalia geography
still preclude event-general operational claims.
